\begin{casebox}
\textbf{知识点小案例：四个线程递增一个计数器为什么扩展性变差。}
特例很小：四个 worker 每处理一条记录都执行一次全局 \code{atomic\_fetch\_add}。内存序即使是 relaxed，目标 cache line 仍然要在多个核心之间转移所有权；如果线程跨 NUMA 节点，转移路径更长。把它改成 worker-local partial 后，热路径只写本地字段，最后再合并，计数语义没变，写共享线的频率却大幅下降。由此推出规律：NUMA 和一致性不是只有大对象才相关，单个热点 cache line 也能成为全局瓶颈。工程上先画谁写哪条 line，再决定全局 atomic、per-worker partial 还是 per-NUMA 合并。
\end{casebox}
